\BourbakiExercisesSection

\begin{enumerate}
\def\labelenumi{\arabic{enumi}.}
\tightlist
\item
  Determine all the group structures on a set of \(n\) elements for \(2 \leqslant n \leqslant 6\) (cf.~§2, Exercise 5). Determine the subgroups and quotient groups of these groups and also their Jordan-Hölder series.
\end{enumerate}

¶ 2. (a) An associative law \((x,y)\mapsto xy\) on a set E is a group law if there exists \(e\in\mathrm{E}\) such that, for all \(x\in\mathrm{E}\) , \(ex=x\) and if, for all \(x\in\mathrm{E}\) , there exists \(x'\) such that \(x^{\prime}x=e\) (show that \(xx^{\prime}=e\) by considering the composition \(x^{\prime}xx^{\prime}\) ; deduce that \(e\) is the identity element).

\begin{enumerate}
\def\labelenumi{(\alph{enumi})}
\setcounter{enumi}{1}
\tightlist
\item
  Show that the same is true if, for all \(x \in E\) , the left translation \(\gamma_{x}\) is a mapping of E onto E and if there exists some \(a \in E\) such that the right translation \(\delta_{a}\) is a mapping of E onto E.
\end{enumerate}

\begin{enumerate}
\def\labelenumi{\arabic{enumi}.}
\setcounter{enumi}{2}
\item
  In a group G every finite non-empty stable subset H is a subgroup of G (cf.~§ 2, Exercise 6).
\item
  Let A and B be two subgroups of a group G.
\end{enumerate}

\begin{enumerate}
\def\labelenumi{(\alph{enumi})}
\item
  Show that the least subgroup containing A and B (that is the subgroup generated by \(\mathbf{A} \cup \mathbf{B}\) ) is identical with the set of compositions of the sequences \((x_{i})_{1 \leqslant i \leqslant 2n + 1}\) of an (arbitrary) odd number of elements such that \(x_{i} \in \mathbf{A}\) for \(i\) odd and \(x_{i} \in \mathbf{B}\) for \(i\) even.
\item
  For AB to be a subgroup of G (in which case it is the subgroup generated by A ∪ B), it is necessary and sufficient that A and B be permutable, that is that AB = BA.
\item
  If A and B are permutable and C is a subgroup containing A, A is permutable with \(B \cap C\) and \(\mathbf{A}(\mathbf{B} \cap \mathbf{C}) = \mathbf{C} \cap (\mathbf{AB})\) .
\end{enumerate}

\begin{enumerate}
\def\labelenumi{\arabic{enumi}.}
\setcounter{enumi}{4}
\item
  If a subgroup of a group G has index 2 it is normal in G.
\item
  Let \((\mathbf{G}_{\alpha})\) be a family of normal subgroups of a group \(\mathbf{G}\) such that \(\bigcap_{\alpha} \mathbf{G}_{\alpha} = \{e\}\) ; show that \(\mathbf{G}\) is isomorphic to a subgroup of the product group \(\prod_{\alpha} (\mathbf{G} / \mathbf{G}_{\alpha})\) .
\item
  If G is the direct product of two subgroups A and B and H is a subgroup of G such that \(A \subset H\) , H is the direct product of A and \(H \cap B\) .
\item
  Let A and A' be two groups and G a subgroup of A × A'. We write:
\end{enumerate}

\[
\mathrm{N} = \mathrm{G} \cap (\mathrm{A} \times \{e \}), \quad \mathrm{H} = \operatorname{pr} _ {1} (\mathrm{G})
\]

\[
\mathrm{N} ^ {\prime} = \mathrm{G} \cap (\{e \} \times \mathrm{A} ^ {\prime}), \quad \mathrm{H} ^ {\prime} = \operatorname{pr} _ {2} (\mathrm{G}).
\]

\begin{enumerate}
\def\labelenumi{(\alph{enumi})}
\tightlist
\item
  Show that N is normal in H and N' is normal in H'; define isomorphisms
\end{enumerate}

\[
\mathrm{H} / \mathrm{N} \rightarrow \mathrm{G} / (\mathrm{N} \times \mathrm{N} ^ {\prime}) \rightarrow \mathrm{H} ^ {\prime} / \mathrm{N} ^ {\prime}.
\]

\begin{enumerate}
\def\labelenumi{(\alph{enumi})}
\setcounter{enumi}{1}
\tightlist
\item
  Suppose that H = A, \(H' = A'\) , that the groups A and \(A'\) are of finite length and that no quotient of a Jordan-Hölder series of A is isomorphic to a quotient of a Jordan-Hölder series of \(A'\) . Show that \(G = A \times A'\) .
\end{enumerate}

\begin{enumerate}
\def\labelenumi{\arabic{enumi}.}
\setcounter{enumi}{8}
\item
  Let G be a group which does not consist only of the identity element. Suppose that there exists a finite subset S of G such that G is generated by S (resp. by the elements \(gsg^{-1}, g \in G, s \in S\) ). Let N be the set of subgroups of G (resp. of normal subgroups of G) distinct from G. Show that N, ordered by inclusion, is inductive. Deduce the existence of a normal subgroup H of G such that G/H is simple.
\item
  Let H be a normal subgroup of a group G contained in the centre of G. Show that if G/H is a monogenous group G is commutative.
\item
  If all the elements of a group G other than the identity element are of order 2, G is commutative; if G is finite, its order n is a power of 2 (argue by induction on n).
\item
  Let G be a group such that, for a fixed integer \(n > 1\) , \((xy)^n = x^n y^n\) for all \(x \in G\) , \(y \in G\) . If \(G^{(n)}\) denotes the set of \(x^n\) , where \(x\) runs through G, and \(G_{(n)}\) the set of \(x \in G\) such that \(x^n = e\) , show that \(G^{(n)}\) and \(G_{(n)}\) are normal subgroups of G; if G is finite, the order of \(G^{(n)}\) is equal to the index of \(G_{(n)}\) . Show that, for all \(x, y\) in G, also \(x^{1-n}y^{1-n} = (xy)^{1-n}\) and deduce that \(x^{n-1}y^n = y^n x^{n-1}\) ; conclude from this that the set of elements of G of the form \(x^{n(n-1)}\) generates a commutative subgroup of G.
\item
  Let G be a group. If S is a simple group, S is said to occur in G if there exist two subgroups H, H' of G with H a normal subgroup of H' such that H'/H is isomorphic to S. Let in(G) denote the set of isomorphism classes of simple groups occurring in G.
\end{enumerate}

\begin{enumerate}
\def\labelenumi{(\alph{enumi})}
\item
  Show that in(G) = ∅ ⇔ G = \{e\}.
\item
  If \(\mathbf{H}\) is a subgroup of \(\mathbf{G}\) , show that \(\operatorname{in}(\mathbf{H}) \subset \operatorname{in}(\mathbf{G})\) ; if moreover \(\mathbf{H}\) is normal, then
\end{enumerate}

\[
\operatorname{in} (\mathrm{G}) = \operatorname{in} (\mathrm{H}) \cup \operatorname{in} (\mathrm{G} / \mathrm{H}).
\]

\begin{enumerate}
\def\labelenumi{(\alph{enumi})}
\setcounter{enumi}{2}
\item
  Let \(\mathbf{G}_1\) and \(\mathbf{G}_2\) be two groups. Show the equivalence of the two following properties:
\item
  \(\operatorname{in}(\mathbf{G}_1) \cap \operatorname{in}(\mathbf{G}_2) = \varnothing\) ;
\end{enumerate}

\begin{enumerate}
\def\labelenumi{(\roman{enumi})}
\setcounter{enumi}{1}
\tightlist
\item
  Every subgroup of \(\mathbf{G}_1\times \mathbf{G}_2\) is of the form \(\mathbf{H}_1\times \mathbf{H}_2\) with \(\mathbf{H}_1\subset \mathbf{G}_1\) and \(\mathbf{H}_2\subset \mathbf{G}_2\) .
\end{enumerate}

(Use Exercise 8.)

*When \(\mathbf{G}_1\) and \(\mathbf{G}_2\) are finite groups, show that (i) and (ii) are equivalent to:

\begin{enumerate}
\def\labelenumi{(\roman{enumi})}
\setcounter{enumi}{2}
\tightlist
\item
  \(\operatorname{Card}(\mathbf{G}_1)\) and \(\operatorname{Card}(\mathbf{G}_2)\) are relatively prime.*
\end{enumerate}

\begin{enumerate}
\def\labelenumi{\arabic{enumi}.}
\setcounter{enumi}{13}
\tightlist
\item
  Let G be a group and H a subgroup of G. Any subset T of G which meets every left coset mod. H in one and only one point is called a representative system of the left cosets mod. H (cf.~Set Theory, II, § 6, no. 2); this condition is equivalent to the following:
\end{enumerate}

The mapping \((x,y)\mapsto xy\) is a bijection of \(\mathbf{T}\times \mathbf{H}\) onto \(\mathbf{G}\) .

\begin{enumerate}
\def\labelenumi{(\alph{enumi})}
\item
  Let \(\mathbf{T}\) be such a representative system; for all \(g \in \mathbf{G}\) and all \(t \in \mathbf{T}\) , let \(x(g, t) \in \mathbf{T}\) and \(y(g, t) \in \mathbf{H}\) be such that \(gt = x(g, t)y(g, t)\) . Let \(\mathbf{S}\) be a subset of \(\mathbf{G}\) generating \(\mathbf{G}\) . Show that the elements \(y(g, t), g \in \mathbf{S}, t \in \mathbf{T}\) generate \(\mathbf{H}\) . (If \(\mathbf{H}'\) denotes the subgroup of \(\mathbf{H}\) generated by these elements, show that \(\mathbf{T}.\mathbf{H}'\) is stable under the \(\boldsymbol{\gamma}_g, g \in \mathbf{S}\) , hence also under the \(\boldsymbol{\gamma}_g, g \in \mathbf{G}\) , and deduce that \(\mathbf{T}.\mathbf{H}' = \mathbf{G}\) , whence \(\mathbf{H}' = \mathbf{H}\) .)
\item
  Suppose that (G:H) is finite. Show that G can be generated by a finite subset if and only if H can.
\end{enumerate}

¶ 15. Let F be a set of stable subgroups of a group with operators G; F is said to satisfy the maximal condition (resp. minimal condition) if every subset of F, ordered by inclusion, has a maximal (resp. minimal) element.

Suppose that the set of all stable subgroups of a group G satisfies the minimal condition.

\begin{enumerate}
\def\labelenumi{(\alph{enumi})}
\item
  Prove that there exists no stable subgroup of G isomorphic to G and distinct from G (argue by reductio ad absurdum by showing that the hypothesis would imply the existence of a strictly decreasing infinite sequence of stable subgroups of G).
\item
  The minimal elements of the set of normal stable subgroups of G not consisting only of e are called minimal normal stable subgroups. Let M be a set of minimal normal stable subgroups of G and S the smallest stable subgroup of G containing all the subgroups belonging to M; show that S is the direct product of a finite number of minimal normal stable subgroups of G (let \((\mathbf{M}_{n})\) be a sequence of minimal normal stable subgroups of G belonging to M and such that \(M_{n+1}\) is not contained in the stable subgroup generated by the union of \(M_{1}, M_{2}, \ldots, M_{n}\) ; let \(S_{k}\) be the stable subgroup generated by the union of the \(M_{n}\) of index \(n \geqslant k\) ; show that \(S_{k+1} = S_{k}\) after a certain rank and therefore that \((\mathbf{M}_{n})\) is a finite sequence; finally use no. 9, Proposition 15).
\item
  If G is a group without operators, show that every minimal normal subgroup M of G is the direct product of a finite number of simple groups isomorphic to one another (let N be a minimal normal subgroup of M; show that M is the smallest subgroup of G containing all the subgroups \(aNa^{-1}\) where a runs through G, and then apply (b) to the group M).
\end{enumerate}

\begin{enumerate}
\def\labelenumi{\arabic{enumi}.}
\setcounter{enumi}{15}
\tightlist
\item
  If the set of stable subgroups of a group with operators G satisfies the maximal and minimal conditions (Exercise 15), G has a Jordan-Hölder series (consider, for a subgroup H of G, a maximal element of the set of normal stable subgroups of H distinct from H).
\end{enumerate}

¶ 17. Let G be a group with operators; a composition series (G \(_{i}\) ) of G is called normal if all the G \(_{i}\) are normal stable subgroups of G; a normal series Σ is called principal if it is strictly decreasing and there exists no normal series distinct from Σ, finer than Σ and strictly decreasing.

\begin{enumerate}
\def\labelenumi{(\alph{enumi})}
\item
  If \((\mathbf{G}_{i})\) and \((\mathbf{H}_{j})\) are two normal series of G, show that there exist two equivalent normal series finer respectively than \((\mathbf{G}_{i})\) and \((\mathbf{H}_{j})\) (apply Schreier's Theorem by considering a suitable domain of operators for G). Give a second proof of this proposition by ``inserting'' the subgroups \(G_{ij}^{\prime} = G_{i} \cap (G_{i+1}H_{j})\) and \(H_{ji}^{\prime} = H_{j} \cap (H_{j+1}G_{i})\) in the sequences \((\mathbf{G}_{i})\) and \((\mathbf{H}_{j})\) respectively.
\item
  If G has a principal series, any two principal series of G are equivalent; for every strictly decreasing normal series \(\Sigma\) there exists a principal series finer than \(\Sigma\) . Deduce that, for G to possess a principal series, it is necessary and sufficient that the set of normal stable subgroups of G satisfy the maximal and minimal conditions.
\item
  If G is a group without operators and possesses a principal series and if the set of subgroups of G satisfies the minimal condition, every quotient group
\end{enumerate}

\(G_{i}/G_{i+1}\) is the direct product of a finite number of simple subgroups isomorphic to one another (use Exercise 15).

¶ 18. (a) Let G be a group with operators generated by the union of a family \((\mathrm{H}_{i})_{i\in\mathbf{I}}\) of simple normal stable subgroups of G. Show that there exists a subset J of I such that G is the restricted sum of the family \((\mathrm{H}_{i})_{i\in\mathbf{J}}\) (apply Zorn's Lemma to the set of subsets K of I such that the subgroup generated by the union of the \(H_{i}\) for \(i\in K\) is the restricted sum of this family).

\begin{enumerate}
\def\labelenumi{(\alph{enumi})}
\setcounter{enumi}{1}
\tightlist
\item
  Let A be a normal stable subgroup of G. Show that there exists a subset \(J_{A}\) of I such that G is the restricted sum of A and the \(H_{i}\) for \(i \in J_{A}\) . Deduce that A is isomorphic to the restricted sum of a subfamily of the \(H_{i}\) .
\end{enumerate}

¶ 19. (a) Let G be a group such that every normal subgroup of G distinct from G is contained in a direct factor of G distinct from G. Show that G is the restricted sum of a family of simple subgroups. (Let K be the subgroup of G generated by the union of all the simple subgroups of G. Suppose K ≠ G; if x \ensuremath{\in} G = K, consider a normal subgroup M of G which is maximal amongst those containing K and not containing x. If G = M' × S, where M' contains M and S ≠ \{e\}, show that x ∉ M' and therefore M' = M; on the other hand, if N is a normal subgroup of S, show that M × N = G and conclude that S is simple, whence a contradiction. Conclude with the aid of Exercise 18.) Obtain the converse (cf.~Exercise 18).

\begin{enumerate}
\def\labelenumi{(\alph{enumi})}
\setcounter{enumi}{1}
\tightlist
\item
  In order that there exist in a group G a family \((\mathrm{N}_{\alpha})\) of normal simple subgroups such that the \(G/N_{\alpha}\) are simple and \(\bigcap_{\alpha} N_{\alpha} = \{e\}\) , it is necessary and sufficient that, for every normal subgroup \(N \neq \{e\}\) of G, there exist a normal subgroup \(N' \neq G\) such that \(NN' = G\) . (To see that the condition is necessary, consider an \(x \in N\) distinct from e and an \(N_{\alpha}\) not containing x. To see that the condition is sufficient, for all \(x \neq e\) in G, consider the normal subgroup N generated by x and a normal subgroup \(N' \neq G\) such that \(NN' = G\) . Show that if a normal subgroup M of G is maximal amongst the normal subgroups of G containing \(N'\) and not containing x, it is maximal in the set of all normal subgroups \(\neq G\) and deduce the conclusion.)
\end{enumerate}

\begin{enumerate}
\def\labelenumi{\arabic{enumi}.}
\setcounter{enumi}{19}
\tightlist
\item
  \begin{enumerate}
  \def\labelenumii{(\alph{enumii})}
  \tightlist
  \item
    Let G be a group, the direct product of two subgroups A and B. Let \(C_{A}\) be the centre of A and let N be a normal subgroup of G such that \(N \cap A = \{1\}\) . Show that \(N \subset C_{A} \times B\) .
  \end{enumerate}
\end{enumerate}

\begin{enumerate}
\def\labelenumi{(\alph{enumi})}
\setcounter{enumi}{1}
\tightlist
\item
  Let \((\mathbf{S}_i)_{i\in I}\) be a finite family of non-commutative simple groups. Show that every normal subgroup of \(\prod_{i\in I}S_i\) is equal to one of the partial products \(\prod_{i\in J}S_i\) where \(J\) is a subset of \(I\) .
\end{enumerate}

\begin{enumerate}
\def\labelenumi{\arabic{enumi}.}
\setcounter{enumi}{20}
\tightlist
\item
  Let G be a group of finite length. Show the equivalence of the following properties:
\end{enumerate}

\begin{enumerate}
\def\labelenumi{(\alph{enumi})}
\item
  G is a product of simple groups isomorphic to one another.
\item
  No subgroup of G distinct from \(\{1\}\) and G is stable under all the automorphisms of G.
\end{enumerate}

¶ 22. Let E be a quasi-group (§ 3, Exercise 6) written multiplicatively, possessing an idempotent e and such that \((xy)(zt) = (xz)(yt)\) for all x, y, z, t. Let \(u(x)\) denote the element of E such that \(u(x)e = x\) and \(v(x)\) the element of E such that \(ev(x) = x\) ; show that the law of composition \((x, y) \mapsto u(x)v(y)\) is a commutative group law on E under which e is identity element, that the mappings \(x \mapsto xe\) and \(y \mapsto ey\) are permutable endomorphisms of this group structure and that \(xy = u(xe)v(ey)\) (start by establishing the identities \(e(xy) = (ex)(ey)\) , \((xy)e = (xe)(ye)\) , \(e(xe) = (ex)e\) ; then note that the relations x = y, ex = ey and xe = ye are equivalent). Obtain the converse.

¶ 23. Consider on a set E an internal law, not everywhere defined, written multiplicatively and satisfying the following conditions:

\begin{enumerate}
\def\labelenumi{(\arabic{enumi})}
\item
  if one of the compositions \((xy)z, x(yz)\) is defined, so is the other and they are equal;
\item
  if \(x, x', y\) are such that \(xy\) and \(x'y\) (resp. \(yx\) and \(yx'\) ) are defined and equal, then \(x = x'\) ;
\item
  for all \(x \in \mathbf{E}\) , there exist three elements \(e_x, e_x'\) and \(x^{-1}\) such that \(e_xx = x\) , \(xe_x' = x\) , \(x^{-1}x = e_x'\) ; \(e_x\) is called the left unit of \(x\) , \(e_x'\) the right unit of \(x\) , \(x^{-1}\) the inverse of \(x\) (by an abuse of language).
\end{enumerate}

\begin{enumerate}
\def\labelenumi{(\alph{enumi})}
\item
  Show that the compositions \(xx^{-1}, x^{-1}e_x, e_x'x^{-1}, e_x e_x, e_x' e_x'\) are defined and that \(xx^{-1} = e_x, x^{-1}e_x = e_x'x^{-1} = x^{-1}, e_x e_x = e_x, e_x' e_x' = e_x'\) .
\item
  Every idempotent \(e\) of \(\mathbf{E}\) (§ 1, no. 4) is a left unit for all the \(x\) such that \(ex\) is defined and right unit for all the \(y\) such that \(ye\) is defined.
\item
  For the composition \(xy\) to be defined, it is necessary and sufficient that the right unit of \(x\) be the same as the left unit of \(y\) (to see that the condition is sufficient, use the relation \(e_y = yy^{-1}\) ); if \(xy = z\) , then \(x^{-1}z = y\) , \(zy^{-1} = x\) , \(y^{-1}x^{-1} = z^{-1}\) , \(z^{-1}x = y^{-1}\) , \(yz^{-1} = x^{-1}\) (the compositions appearing on the left-hand sides of these relations being defined).
\item
  For any two idempotents \(e, e'\) of \(\mathbf{E}\) , let \(\mathbf{G}_{e, e'}\) denote the set of \(x \in \mathbf{E}\) such that \(e\) is left unit and \(e'\) is right unit of \(x\) . Show that, if \(\mathbf{G}_{e, e}\) is non-empty, it is a group under the law induced by that on \(\mathbf{E}\) .
\end{enumerate}

E is called a groupoid if it also satisfies the following condition:

\begin{enumerate}
\def\labelenumi{(\arabic{enumi})}
\setcounter{enumi}{3}
\tightlist
\item
  for all idempotents \(e, e', G_{e, e'}\) is non-empty.
\end{enumerate}

\begin{enumerate}
\def\labelenumi{(\alph{enumi})}
\setcounter{enumi}{4}
\item
  In a groupoid E, if \(a \in G_{e,e'}\) , show that \(x \mapsto xa\) is a bijective mapping of \(G_{e,e}\) onto \(G_{e,e'}, y \mapsto ay\) a bijective mapping of \(G_{e',e'}\) onto \(G_{e,e'}\) and \(x \mapsto a^{-1}xa\) an isomorphism of the group \(G_{e,e}\) onto the group \(G_{e',e'}\) .
\item
  Show that the law defined in § 1, Exercise 4(b) determines a groupoid structure on the set \(\Psi\) , if all the sets of the family \(\mathfrak{F}\) are equipotent.
\end{enumerate}

\begin{enumerate}
\def\labelenumi{\arabic{enumi}.}
\setcounter{enumi}{23}
\tightlist
\item
  \begin{enumerate}
  \def\labelenumii{(\alph{enumii})}
  \tightlist
  \item
    Given any set E, consider in the set \(E \times E\) the law of composition, not everywhere defined, written multiplicatively, under which the composition of \((x, y)\) and \((y', z)\) is only defined if \(y' = y\) and in this case has the value \((x, z)\) . Show that \(\mathbf{E} \times \mathbf{E}\) , with this law of composition, is a groupoid (Exercise 23).
  \end{enumerate}
\end{enumerate}

\begin{enumerate}
\def\labelenumi{(\alph{enumi})}
\setcounter{enumi}{1}
\tightlist
\item
  Let \((x, y) \equiv (x', y')\) (R) be an equivalence relation compatible with the above composition (that is such that \((x, y) \equiv (x', y')\) and \((y, z) \equiv (y', z')\) imply \((x, z) \equiv (x', z'))\) ; suppose further that the relation R satisfies the following condition: for all \(x, y, z\) there exists one and only one \(t \in \mathbf{E}\) such that \((x, y) \equiv (z, t)\) and there exists a \(u \in \mathbf{E}\) such that \((x, y) \equiv (u, z)\) .
\end{enumerate}

Show that under these conditions the quotient structure by R of the groupoid structure on \(E \times E\) is a group structure (prove first that the quotient law is everywhere defined, since, if \(\dot{x}, \dot{y}, \dot{z}\) are three classes such that \(\dot{x}\dot{y} = \dot{x}\dot{z}\) , then \(\dot{y} = \dot{z}\) ; establish finally that, for all \(x \in E, y \in E, (x, x) \equiv (y, y)\) ).

\begin{enumerate}
\def\labelenumi{(\alph{enumi})}
\setcounter{enumi}{2}
\tightlist
\item
  Let G be the group obtained by giving the quotient of \(E \times E\) by R the above structure. Let a be any element of E; if, for all \(x \in E\) , \(f_{a}(x)\) denotes the class mod. R of \((a, x)\) , show that \(f_{a}\) is a bijective mapping of E onto G and that the relation \((x, y) \equiv (x', y')\) is equivalent to the relation
\end{enumerate}

\[
f _ {a} (x) \left(f _ {a} \left(x ^ {\prime}\right)\right) ^ {- 1} = f _ {a} (y) \left(f _ {a} \left(y ^ {\prime}\right)\right) ^ {- 1}.
\]

For G to be commutative, it is necessary and sufficient that the relation \((x,y)\equiv (x',y')\) imply \((x,x')\equiv (y,y')\) .

¶ 25. Let E be a set and f a mapping of \(E^{m}\) into E; we write

\[
f (x _ {1}, x _ {2}, \dots , x _ {m}) = x _ {1} x _ {2} \dots x _ {m};
\]

suppose that \(f\) satisfies the following conditions:

\[
(x _ {1} x _ {2} \dots x _ {m}) x _ {m + 1} \dots x _ {2 m - 1} = x _ {1} (x _ {2} x _ {3} \dots x _ {m + 1}) x _ {m + 2} \dots x _ {2 m - 1}\tag{1}
\]

identically;

\begin{enumerate}
\def\labelenumi{(\arabic{enumi})}
\setcounter{enumi}{1}
\tightlist
\item
  For all \(a_1, a_2, \ldots, a_{m-1}\) , the mappings
\end{enumerate}

\[
x \mapsto x a _ {1} a _ {2} \dots a _ {m - 1}
\]

\[
x \mapsto a _ {1} a _ {2} \dots a _ {m - 1} x
\]

are bijections of E onto E.

\begin{enumerate}
\def\labelenumi{(\alph{enumi})}
\tightlist
\item
  Show that identically
\end{enumerate}

\[
(x _ {1} x _ {2} \dots x _ {m}) x _ {m + 1} \dots x _ {2 m - 1} = x _ {1} x _ {2} \dots x _ {i} (x _ {i + 1} \dots x _ {i + m}) x _ {i + m + 1} \dots x _ {2 m - 1}
\]

for every index \(i\) such that \(1 \leqslant i \leqslant m - 1\) (argue by induction on \(i\) , by considering the element

\[
\left(\left(x _ {1} x _ {2} \dots x _ {m}\right) x _ {m + 1} \dots x _ {2 m - 1}\right) a _ {1} a _ {2} \dots a _ {m - 1}).
\]

\begin{enumerate}
\def\labelenumi{(\alph{enumi})}
\setcounter{enumi}{1}
\tightlist
\item
  For all \(a_1, a_2, \ldots, a_{m-2}\) , there exists \(u \in E\) such that for all \(x\)
\end{enumerate}

\[
x = x a _ {1} a _ {2} \dots a _ {m - 2} u = u a _ {1} a _ {2} \dots a _ {m - 2} x.
\]

\begin{enumerate}
\def\labelenumi{(\alph{enumi})}
\setcounter{enumi}{2}
\tightlist
\item
  In the set \(\mathbf{E}^k\) of sequences \((u_1, u_2, \ldots, u_k)\) of \(k\) elements of \(\mathbf{E}\) ( \(1 \leqslant k \leqslant m - 1\) ) consider the equivalence relation \(\mathbf{R}_k\) which states: for all \(x_1, x_2, \ldots, x_{m-k}, u_1 u_2 \ldots u_k x_1 x_2 \ldots x_{m-k} = v_1 v_2 \ldots v_k x_1 x_2 \ldots x_{m-k}\) ; let \(\mathbf{E}_k\) denote the quotient set \(\mathbf{E}^k / \mathbf{R}_k\) and \(\mathbf{G}\) the sum set (Set Theory, II, §4, no. 8) of \(\mathbf{E}_1 = \mathbf{E}, \mathbf{E}_2, \ldots, \mathbf{E}_{m-1}\) . Let \(\alpha \in \mathbf{E}_i\) , \(\beta \in \mathbf{E}_j\) ; if \((u_1, u_2, \ldots, u_i)\) is a sequence of class \(\alpha\) , \((v_1, \ldots, v_j)\) a sequence of class \(\beta\) , consider the sequence
\end{enumerate}

\[
\left(u _ {1}, u _ {2}, \dots , u _ {i}, v _ {1}, \dots , v _ {j}\right)
\]

in \(\mathbf{E}^{i + j}\) if \(i + j < m\) , the sequence

\[
\left(u _ {1}, u _ {2}, \dots , u _ {i + j - m}, \left(u _ {i + j - m + 1} \dots u _ {i} v _ {1} v _ {2} \dots v _ {j}\right)\right)
\]

in \(E^{i+j-m+1}\) if \(i+j\geqslant m\) ; show that the class of this sequence in \(E_{i+j}\) (resp. \(E_{i+j-m+1}\) ) depends only on the classes \(\alpha\) and \(\beta\) ; it is denoted by \(\alpha.\beta\) . Show that a group law is thus defined on G, that \(H=E_{m-1}\) is a normal subgroup of G and that the quotient group G/H is cyclic of order m-1; prove finally that E is identical with a class (mod. H) which generates G/H and that \(x_{1}x_{2}\ldots x_{m}\) is just the composition, in the group G, of the sequence \((x_{1},x_{2},\ldots,x_{m})\) .

¶ 26. Let G, G' be two groups, \(f: G \to G'\) a mapping such that, for two arbitrary elements \(x, y\) of G, \(f(xy) = f(x)f(y)\) or \(f(xy) = f(y)f(x)\) . It is proposed to prove that \(f\) is a homomorphism of G into G' or a homomorphism of G into the opposite group G'0 (in other words, either \(f(xy) = f(x)f(y)\) for every ordered pair \((x, y)\) or \(f(xy) = f(y)f(x)\) for every ordered pair \((x, y)\) ).

\begin{enumerate}
\def\labelenumi{(\alph{enumi})}
\item
  Show that the set \(N = f(e')\) ( \(e'\) the identity element of \(G'\) ) is a normal subgroup of G and that f factors into \(G \xrightarrow{p} G/N \xrightarrow{g} G'\) , where g is injective. Attention may therefore be confined to the case where f is injective, which will be assumed throughout what follows.
\item
  Show that if \(xy = yx\) then \(f(x)f(y) = f(y)f(x)\) (consider \(f(x^2y)\) and \(f(x^2y^2)\) , expressing them in several ways).
\item
  Show that if \(f(xy) = f(x)f(y)\) then \(f(yx) = f(y)f(x)\) . (Show with the aid of (b) that attention may be confined to the case where \(xy \neq yx\) .)
\item
  Show that \(f(xyx) = \hat{f}(x)f(y)f(x)\) . (Use (a) or (b) according to whether \(xy = yx\) or \(xy \neq yx\) ; in the second case, consider \(f(x^2y)\) and \(f(yx^2)\) .)
\item
  Let A be the set of \(x \in G\) such that \(f(xy) = f(x)f(y)\) for all \(y \in G\) and B the set of \(x \in G\) such that \(f(xy) = f(y)f(x)\) for all \(y \in G\) . Show that \(A \neq G\) , \(B \neq G\) and \(A \cup B = G\) is an impossible situation. (By virtue of (b) there would then exist \(x, y\) in A such that \(f(xy) = f(x)f(y) \neq f(y)f(x)\) and \(u, v\) in B such that \(f(uv) = f(v)f(u) \neq f(u)f(v)\) . Deduce from this a contradiction by considering successively the two possibilities \(xu \in A\) , \(xu \in B\) ; in the first case consider \(f(xuv)\) and in the second \(f(yxu)\) .)
\end{enumerate}

The rest of the proof consists of proving that \(\mathbf{A} \cup \mathbf{B} \neq \mathbf{G}\) is impossible, arguing by reductio ad absurdum. If \(\mathbf{A} \cup \mathbf{B} \neq \mathbf{G}\) , there exist \(a, b, c\) in \(\mathbf{G}\) such that

\[
f (a b) = f (a) f (b) \neq f (b) f (a) \quad \text { and } \quad f (a c) = f (c) f (a) \neq f (a) f (c).
\]

\begin{enumerate}
\def\labelenumi{(\alph{enumi})}
\setcounter{enumi}{5}
\item
  Show that \(f(c)f(a)f(b) = f(b)f(a)f(c)\) . (Consider two cases according to whether \(bc \neq cb\) or \(bc = cb\) . In the first case, consider \(f(bac)\) ; in the second, use (b) and (c) and consider \(f(abc), f(bca)\) and \(f(bac)\) .)
\item
  Consider \(f(abac)\) and obtain a contradiction.
\end{enumerate}

